\documentclass[a4paper,10pt]{article}
\usepackage[francais]{babel}
\usepackage{url}

\selectlanguage{francais}
../letter-ensiie.sty

\selectlanguage{francais}
\enCouleur

\begin{document}


\name{St\'ephane \textsc{Genaud}}
\fonction{Professor at ENSIIE}
\telephoneillkirch{+33 (0)3688 54542}
\email{genaud@ensiie.fr}

\begin{letter}[Recommendation de Magdeleine \textsc{Lebrun}]{à}%
{au responsable du MAE Ingénieur manager, EM Strasbourg}

Chers Collègues,

Je connais Magdeleine Lebrun depuis sa  première année de scolarité dans l'école
d'ingénieur en  informatique ENSIIE  alors que je  dirigeais l'antenne  locale à
Strasbourg,  et  dans laquelle  j'interviens  également  comme enseignant  et  y
supervise les stages.  Je donne des éléments d'appréciation de la candidature de
Magdeleine à ces deux titres.


La formation d'ingénieur de l'ENSIIE est une formation en informatique et
mathématiques qui est dispensée à la fois sur le site
historique principal, à Evry, et à son antenne de Strasbourg selon les modalités
suivantes: un tronc commun se déroule sur deux semestres, avec des contenus 
pédagogiques et des examens communs aux deux sites, puis les semestres 3, 4 et 5
sont constitués à 80\% d'options axées sur les spécialités recherche des sites.
Le dernier semestre est un stage long de type professionnel ou recherche. 

Tout  au  long de  sa  scolarité,  Magdeleine a  obtenu  de  très bons  résulats
scolaires, la plaçant dans le premier tiers de la promotion. Elle s'est montrée
assidue et rigoureuse dans son travail de manière constante, avec un très grand
sens des responsabilités. 
 

Je  retiens  également  son  grand  sens   de  la  responsabilité  dans  la  vie
para-scolaire.  Elle  a été  élue présidente  du Bureau des  Elèves (BDE)  de sa
promotion.  Pour avoir travaillé étroitement avec le BDE sur de nombreux projets
impliquant la  direction de l'école,  je peux  témoigner de son  souci permanent
d'organiser ces actions  collectives de manière rigoureuse pour  s'en assurer la
réussite.

Ayant  également suivi  les  choix  de stage  des  étudiants,  j'ai remarqué  la
curiosité de Magdeleine qui a varié les environnements choisis pour ses stages :
une action  humanitaire au Laos en  première année, une startup  très pointue en
informatique  en deuxième  année, et  une  grande entreprise  dans les  services
numériques (ESN) en fin de cursus. Elle a donc pu se rendre compte des différences
notoires, particulièrement en terme d'organisation du travail, à travers ces
expériences variées. Cette dernière expérience professionnelle en ESN étant
effectuée sous la forme d'un apprentissage, elle a pu mûrir une vision très
claire de son futur projet professionnel.


Je suis persuadée que Magdeleine a de très grandes qualités humaines et
d'organisation, alliées à des compétences techniques et scientifiques en
informatique, qui lui permettront sans nul doute de mener une brillante
carrière dans la gestion de projet.

Pour m'être entretenu  avec elle de son projet professionnel,  je suis entièrement
convaincu par  sa motivation  à acquérir les  compétences transversales  qui lui
manquent et ainsi augmenter sa polyvalence.  Magdeleine souhaitait faire acte de
candidature au MAE AlsaceTech auquel les étudiants de l'ENSIIE ne sont plus
éligibles depuis la fermeture de l'antenne strasbourgeoise. Elle reste néanmoins
persuadée de la pertinence de cette formation pour elle et pour ces raisons, je recommande au
plus haut point sa candidature au  MAE Ingénieur manager de l'EM Strasbourg.


Je vous prie, Madame, Monsieur, d'agréer mes sentiments respectueux.

\end{letter}
\begin{flushright}
\includegraphics[width=.2\textwidth]{signgenaud.jpg}
\end{flushright}
\end{document}
